\section{Personal Pilot Integration in the Unified Mobile Application}

Stage~5 begins by connecting the new phone shell to Pilot's existing personal authorities.  It
does not create a replacement identity store, task database or household model.  The
\texttt{PersonalPilotProjection} is a narrow, read-only mobile adapter over Production Private
Identity and Pilot Inbox.

\subsection{Phone-bound access}

The Personal Pilot request is sent to the mother over trusted HTTPS.  It contains the exact
purpose and persona and is signed by the phone's non-extractable Ed25519 possession key.  The
mother verifies its own signed phone certificate, the current capability lease, the
\texttt{inbox.read} scope, the authoritative trusted-device record, the persona binding and the
signature over the complete request.  A copied certificate without the phone key, a revoked
device, an expired lease, a modified request or a request for another persona fails closed.

\subsection{Identity and guardian boundary}

The resulting projection includes the signed-in person's active profile.  An adult also sees
profiles for children whose explicit guardian identifier names that adult.  It does not return
another adult's profile through this private view, and it does not disclose another person's
tasks, reminders, contacts or messages.  Recovery-code hashes and device public keys remain
outside the projection.  Loading the mobile view does not activate that person on a shared
television or grant shared-screen authority.

\subsection{Real personal work rather than fixture cards}

Lists, tasks, delegations and acceptance states are read directly from Pilot Inbox.  The same
authority remains responsible for private versus household list scope, recipient acceptance,
snooze, completion and provider receipts.  The phone renders a compact summary and recent real
tasks after pairing; the established Unified Inbox controls continue to handle recipient task
acceptance and lifecycle changes.  Fixture cards remain only when no mother brain is connected.

The projection also carries bounded reminder and contact counts so the interface can truthfully
indicate available personal data while their complete Stage~5 interaction surfaces remain open.
It does not claim reminder, calendar, contact, service, file, device, memory or Guardian
integration is complete merely because those records can be counted.

\subsection{Verification and remaining work}

Automated evidence creates two independent adults plus one child controlled by a guardian, binds a
phone to one adult, and verifies that only that adult, their child and their own tasks, list,
reminder summary and contact count are returned.  Altered persona requests and invalid phone
signatures are rejected.  The optimized production phone application also type-checks and builds
its unified mobile route successfully.

\subsection{Governed reminder lifecycle}

The phone can now schedule a reminder against one of the active person's real tasks.  Exact-time
and closing-time reminders accept a future local date and time.  They may run once, daily, on
weekdays or weekly.  Completing a repeating timed reminder closes the current occurrence and
atomically creates exactly one next occurrence; duplicate signed requests return the original
result.  A person can also snooze a reminder for a future time, complete it or cancel it.

Arrival, departure and journey triggers require more authority.  The mobile request must name an
active Pilot Location consent containing the reminder-location scope; that consent must
belong to the same persona and exact phone device that signs the request.  Merely entering a place
name or possessing another household member's permission does not suffice.  Pilot retains the
human-readable location label and permission reference but explicitly records that raw location
history is not retained.  The phone explains when permission is absent instead of silently
downgrading the request to a timed reminder.

Creation and lifecycle changes use canonical phone signatures, the current capability lease and
durable idempotency keys.  Reusing a key with changed content, altering the persona, acting on
another person's reminder, using a past time or providing an unsupported repetition fails closed.
Automated evidence covers exact repeat preservation, single next-occurrence creation, snooze,
location denial without consent, success with same-device consent and cross-person rejection.

This closes the first three coded items in Stage~5: identity, tasks and reminders.  Stage~5
continues with complete calendar, private-contact, communication, shopping, booking, files,
saved-item, device, television, learning, games, memory, Guardian and AION Native surfaces, then real provider
qualification.  Those capabilities must continue to reuse their existing authorities and exact
approval boundaries.
